% Document/LinearAlgebra/VectorSpaces/LinearMaps/Properties.tex
% Section: Injectivity, Surjectivity, Bijectivity Characterizations

\newpage
\section{Injectivity, Surjectivity, and Bijectivity}
\label{sec:inj-surj-bij}

\begin{remark}[Overview]
    For a general function between sets, injectivity, surjectivity, and bijectivity are
    independent properties. For a \emph{linear} map, however, these set-theoretic conditions
    have multiple equivalent reformulations involving kernels, images, dimensions,
    independence, spanning, lattice operations, and categorical properties.

    The three theorems below collect these equivalences into exhaustive lists. Each list is
    proved by reducing every condition to its core case --- a trivial kernel for injectivity,
    a full range for surjectivity --- via a few short implication cycles through that core;
    bijectivity then follows by conjoining the two theorems.
    The key insight running through all three is:
    \begin{itemize}
        \item \textbf{Injectivity} is about \emph{preserving independence}: an injective
              map does not collapse distinct directions. Its hallmark is a trivial kernel.
        \item \textbf{Surjectivity} is about \emph{preserving generation}: a surjective
              map covers the entire target. Its hallmark is that every generating set maps
              to a generating set.
        \item \textbf{Bijectivity} combines both: it is about \emph{preserving bases},
              i.e., both independence and generation simultaneously.
    \end{itemize}
    In finite dimensions, injectivity and surjectivity become equivalent for maps between
    spaces of the same dimension (Theorem~\ref{thm:equal-dim-equiv}) --- a consequence of
    the fact that on finite sets of the same cardinal injective functions are surjective and vice versa.
\end{remark}

\begin{theorem}[Characterizations of Injectivity]\label{thm:injectivity}
    For a linear map $T: U \to V$, the following are equivalent:
    \begin{enumerate}
        \item $\Ker{T} = \{0\}$
        \item $T$ is injective
        \item $\InverseImage{T}{\DirectImage{T}{W}} = W$ for every subspace
              $W \Subspace[K] U$
        \item $W \mapsto \DirectImage{T}{W}$ is injective on the subspace lattice of $U$
        \item $\Dim[K]{\DirectImage{T}{W}} = \Dim[K]{W}$ for every subspace
              $W \Subspace[K] U$
        \item $\Dim[K]{\DirectImage{T}{W}} = \Dim[K]{W}$ for every
              \emph{finite-dimensional} subspace $W \Subspace[K] U$
        \item $\Dim[K]{\DirectImage{T}{W}} = \Dim[K]{W}$ for every
              \emph{one-dimensional} subspace $W \Subspace[K] U$
              (equivalently, $T(v) \neq 0$ for every $v \neq 0$)
        \item $T$ sends every linearly independent set to a linearly independent set
        \item $T$ sends every basis of $U$ to a linearly independent set
        \item $T$ sends some basis of $U$ to a linearly independent set
        \item $T$ is a split monomorphism (has a left inverse):
              $\Exists :{S: V \to U}{\Composition{S}{T} = \Identity{U}}$
        \item $T$ is monic (left-cancellable):
              $\Forall :{S_1, S_2: W \to U}{\Composition{T}{S_1} = \Composition{T}{S_2}
              \implies S_1 = S_2}$
    \end{enumerate}
\end{theorem}

\begin{proof}
    We show every condition is equivalent to the trivial-kernel condition~(1) by proving
    three cycles through it:
    \[
        (1)\Rightarrow(2)\Rightarrow(3)\Rightarrow(4)\Rightarrow(1), \qquad
        (1)\Rightarrow(5)\Rightarrow(6)\Rightarrow(7)\Rightarrow(1), \qquad
        (1)\Rightarrow(8)\Rightarrow(9)\Rightarrow(10)\Rightarrow(11)\Rightarrow(12)\Rightarrow(1).
    \]

    \textbf{(1) $\Rightarrow$ (2)}:
    Let $T(u_1) = T(u_2)$. Then $T(u_1 - u_2) = 0$, so
    $u_1 - u_2 \in \Ker{T} = \{0\}$, hence $u_1 = u_2$.

    \textbf{(2) $\Rightarrow$ (3)}:
    The inclusion $W \subseteq \InverseImage{T}{\DirectImage{T}{W}}$ is trivial.
    Conversely, if $u \in \InverseImage{T}{\DirectImage{T}{W}}$, then $T(u) = T(w)$ for some
    $w \in W$. Injectivity gives $u = w \in W$.

    \textbf{(3) $\Rightarrow$ (4)}:
    If $\DirectImage{T}{W_1} = \DirectImage{T}{W_2}$, then
    $W_1 = \InverseImage{T}{\DirectImage{T}{W_1}}
    = \InverseImage{T}{\DirectImage{T}{W_2}} = W_2$.

    \textbf{(4) $\Rightarrow$ (1)}:
    Let $k \in \Ker{T}$. Then $\DirectImage{T}{\Span[K]{\{k\}}} = \Span[K]{\{T(k)\}}
    = \{0\} = \DirectImage{T}{\{0\}}$, so injectivity of $W \mapsto \DirectImage{T}{W}$
    forces $\Span[K]{\{k\}} = \{0\}$, i.e.\ $k = 0$. Hence $\Ker{T} = \{0\}$.

    \textbf{(1) $\Rightarrow$ (5)}:
    Let $W \Subspace[K] U$ with basis $B$. By Proposition~\ref{prop:span-image},
    $\DirectImage{T}{B}$ generates $\DirectImage{T}{W}$, and it is linearly independent:
    $\sum_{i \in n} a_i T(b_i) = 0$ gives $\sum_{i \in n} a_i b_i \in \Ker{T} = \{0\}$, so all
    $a_i = 0$. Since $\Ker{T} = \{0\}$ also makes $T$ injective on $B$, we have
    $\Card{\DirectImage{T}{B}} = \Card{B}$; thus $\DirectImage{T}{B}$ is a basis of
    $\DirectImage{T}{W}$ and $\Dim[K]{\DirectImage{T}{W}} = \Dim[K]{W}$.

    \textbf{(5) $\Rightarrow$ (6)}:
    Every finite-dimensional subspace is a subspace, so the equality in (5) holds in
    particular for every finite-dimensional $W \Subspace[K] U$.

    \textbf{(6) $\Rightarrow$ (7)}:
    Every one-dimensional subspace is finite-dimensional, so the equality in (6) holds in
    particular for every one-dimensional $W \Subspace[K] U$.

    \textbf{(7) $\Rightarrow$ (1)}:
    Item (7) says $T(v) \neq 0$ for every $v \neq 0$; equivalently, $T(v) = 0$ implies
    $v = 0$. Thus $\Ker{T} = \{0\}$.

    \textbf{(1) $\Rightarrow$ (8)}:
    Let $L \subseteq U$ be linearly independent and suppose
    $\sum_{i \in n} a_i T(\ell_i) = 0$ for distinct $\ell_i \in L$. Then
    $T\bigl(\sum_{i \in n} a_i \ell_i\bigr) = 0$, so $\sum_{i \in n} a_i \ell_i \in \Ker{T}
    = \{0\}$; linear independence of $L$ forces all $a_i = 0$. Hence $\DirectImage{T}{L}$ is
    linearly independent.

    \textbf{(8) $\Rightarrow$ (9)}:
    Every basis is linearly independent.

    \textbf{(9) $\Rightarrow$ (10)}:
    Every basis is some basis.

    \textbf{(10) $\Rightarrow$ (11)}:
    Suppose some basis $B$ is sent to a linearly independent set $\DirectImage{T}{B}$.
    Extend $\DirectImage{T}{B}$ to a basis $B' = \DirectImage{T}{B} \Union C$ for $V$
    (Theorem~\ref{thm:basis-extension}).
    Define $S: V \to U$ by $S(T(b)) = b$ for $b \in B$ and $S(c) = 0$ for $c \in C$,
    extended linearly (Theorem~\ref{thm:vector-space-universal}).
    Then $(\Composition{S}{T})(b) = b$ for all $b \in B$, so
    $\Composition{S}{T} = \Identity{U}$.

    \textbf{(11) $\Rightarrow$ (12)}:
    Suppose $\Composition{S}{T} = \Identity{U}$ and
    $\Composition{T}{S_1} = \Composition{T}{S_2}$. Composing with $S$ on the left:
    $S_1 = \Composition{(\Composition{S}{T})}{S_1}
    = \Composition{S}{(\Composition{T}{S_1})}
    = \Composition{S}{(\Composition{T}{S_2})}
    = \Composition{(\Composition{S}{T})}{S_2} = S_2$.

    \textbf{(12) $\Rightarrow$ (1)}:
    Let $\iota: \Ker{T} \hookrightarrow U$ be the inclusion and $0: \Ker{T} \to U$ the zero
    map. For any $u \in \Ker{T}$: $T(\iota(u)) = T(u) = 0 = T(0(u))$, so
    $\Composition{T}{\iota} = \Composition{T}{0}$. By the monic property,
    $\iota = 0$, hence $\Ker{T} = \{0\}$.
\end{proof}

\begin{remark}[Split monomorphism and monomorphism, diagrammatically]
    The two categorical characterizations have the following pictures. Condition~(11), that
    $T$ is a \textbf{split monomorphism}, asserts a left inverse (retraction) $S$ making the
    triangle commute:
    \begin{center}
    \begin{tikzcd}
        U \arrow[r, "T"] \arrow[dr, "\Identity{U}"'] & V \arrow[d, "S"] \\
        & U
    \end{tikzcd}
    \end{center}
    Condition~(12), that $T$ is a \textbf{monomorphism}, is left-cancellability of $T$ against
    any parallel pair $S_1, S_2: W \to U$:
    \begin{center}
    \begin{tikzcd}
        W \arrow[r, shift left, "S_1"] \arrow[r, shift right, "S_2"'] & U \arrow[r, "T"] & V
    \end{tikzcd}
    \qquad
    $\Composition{T}{S_1} = \Composition{T}{S_2} \;\implies\; S_1 = S_2$.
    \end{center}
    A retraction always makes $T$ monic; over a field the converse holds as well, which is the
    substance of the equivalence (11)$\Leftrightarrow$(12).
\end{remark}

\newpage
\begin{theorem}[Characterizations of Surjectivity]\label{thm:surjectivity}
    For a linear map $T: U \to V$, the following are equivalent:
    \begin{enumerate}
        \item $\Rng{T} = V$
        \item $T$ is surjective
        \item $\DirectImage{T}{\InverseImage{T}{W}} = W$ for every subspace
              $W \Subspace[K] V$
        \item $W \mapsto \InverseImage{T}{W}$ is injective on the subspace lattice of $V$
        \item $\Codim[K]{\InverseImage{T}{W}}{U} = \Codim[K]{W}{V}$ for every subspace
              $W \Subspace[K] V$
        \item $\Codim[K]{\InverseImage{T}{W}}{U} = \Codim[K]{W}{V}$ for every
              \emph{finite-codimensional} subspace $W \Subspace[K] V$
        \item $\Codim[K]{\InverseImage{T}{W}}{U} = \Codim[K]{W}{V}$ for every
              subspace $W \Subspace[K] V$ of \emph{codimension one} (every hyperplane)
        \item $T$ sends every generating system of $U$ to a generating system of $V$
        \item $T$ sends every basis of $U$ to a generating system of $V$
        \item $T$ sends some basis of $U$ to a generating system of $V$
        \item $T$ is a split epimorphism (has a right inverse):
              $\Exists :{S: V \to U}{\Composition{T}{S} = \Identity{V}}$
        \item $T$ is epic (right-cancellable):
              $\Forall :{S_1, S_2: V \to W}{\Composition{S_1}{T} = \Composition{S_2}{T}
              \implies S_1 = S_2}$
    \end{enumerate}
\end{theorem}

\begin{proof}
    Dually, every condition is equivalent to the full-range condition~(1) by three cycles
    through it:
    \[
        (1)\Rightarrow(2)\Rightarrow(3)\Rightarrow(4)\Rightarrow(1), \qquad
        (1)\Rightarrow(5)\Rightarrow(6)\Rightarrow(7)\Rightarrow(1), \qquad
        (1)\Rightarrow(8)\Rightarrow(9)\Rightarrow(10)\Rightarrow(11)\Rightarrow(12)\Rightarrow(1).
    \]

    \textbf{(1) $\Rightarrow$ (2)}:
    This is the definition of surjectivity.

    \textbf{(2) $\Rightarrow$ (3)}:
    By Proposition~\ref{prop:image-preimage-formulas}(2),
    $\DirectImage{T}{\InverseImage{T}{W}} = W \cap \Rng{T} = W \cap V = W$.

    \textbf{(3) $\Rightarrow$ (4)}:
    If $\InverseImage{T}{W_1} = \InverseImage{T}{W_2}$, then
    $W_1 = \DirectImage{T}{\InverseImage{T}{W_1}}
    = \DirectImage{T}{\InverseImage{T}{W_2}} = W_2$.

    \textbf{(4) $\Rightarrow$ (1)}:
    Since $\InverseImage{T}{\Rng{T}} = U = \InverseImage{T}{V}$, injectivity of
    $W \mapsto \InverseImage{T}{W}$ forces $\Rng{T} = V$.

    \textbf{(1) $\Rightarrow$ (5)}:
    Let $W \Subspace[K] V$ and pick a complement $C$ of $\InverseImage{T}{W}$ in $U$.

    \textit{Claim}: $\DirectImage{T}{C}$ is a complement of $W$ in $V$.
    If $T(c) \in W$ for $c \in C$, then $c \in \InverseImage{T}{W} \cap C = \{0\}$.
    For any $v \in V$, $\Rng{T} = V$ gives $v = T(w' + c) = T(w') + T(c)$ with
    $w' \in \InverseImage{T}{W}$, $c \in C$, so $v \in W + \DirectImage{T}{C}$.

    The map $T|_C: C \to \DirectImage{T}{C}$ is injective: if $T(c) = 0$, then
    $c \in \Ker{T} \subseteq \InverseImage{T}{W}$ (since $0 \in W$), so
    $c \in \InverseImage{T}{W} \cap C = \{0\}$. By Theorem~\ref{thm:injectivity}(5),
    $\Dim[K]{C} = \Dim[K]{\DirectImage{T}{C}}$. Thus
    $\Codim[K]{\InverseImage{T}{W}}{U} = \Dim[K]{C}
    = \Dim[K]{\DirectImage{T}{C}} = \Codim[K]{W}{V}$.

    \textbf{(5) $\Rightarrow$ (6)}:
    Every finite-codimensional subspace is a subspace, so the equality in (5) holds in
    particular for every finite-codimensional $W \Subspace[K] V$.

    \textbf{(6) $\Rightarrow$ (7)}:
    Every subspace of codimension one is finite-codimensional, so the equality in (6) holds
    in particular for every codimension-one $W \Subspace[K] V$.

    \textbf{(7) $\Rightarrow$ (1)}:
    Item (7) says every hyperplane $W \Subspace[K] V$ has
    $\Codim[K]{\InverseImage{T}{W}}{U} = 1$, so in particular $\InverseImage{T}{W} \neq U$:
    some $u \in U$ has $T(u) \notin W$, whence $\Rng{T} \not\SubsetEq W$. Thus $\Rng{T}$ is
    contained in no hyperplane of $V$. But every proper subspace of $V$ lies in a hyperplane
    (extend a basis of the subspace to a basis of $V$ by Theorem~\ref{thm:basis-extension}
    and omit one of the added vectors), so $\Rng{T}$ is not proper; that is, $\Rng{T} = V$.

    \textbf{(1) $\Rightarrow$ (8)}:
    Let $G$ generate $U$. By Proposition~\ref{prop:span-image},
    $\Span[K]{\DirectImage{T}{G}} = \DirectImage{T}{\Span[K]{G}} = \Rng{T} = V$, so
    $\DirectImage{T}{G}$ generates $V$.

    \textbf{(8) $\Rightarrow$ (9)}:
    Every basis is a generating system.

    \textbf{(9) $\Rightarrow$ (10)}:
    Every basis is some basis.

    \textbf{(10) $\Rightarrow$ (11)}:
    Suppose some basis $B$ of $U$ maps to a generating system $\DirectImage{T}{B}$ for $V$.
    Let $B_V$ be any basis for $V$. For each $b_V \in B_V$, write
    $b_V = \sum_{b \in F} a_b T(b)$ for some finite $F \subseteq B$.
    Define $S: V \to U$ by $S(b_V) = \sum_{b \in F} a_b b$ and extend linearly.
    Then $(\Composition{T}{S})(b_V) = b_V$ for all $b_V \in B_V$, so
    $\Composition{T}{S} = \Identity{V}$.

    \textbf{(11) $\Rightarrow$ (12)}:
    Suppose $\Composition{T}{S} = \Identity{V}$ and
    $\Composition{S_1}{T} = \Composition{S_2}{T}$. Composing with $S$ on the right:
    $S_1 = \Composition{S_1}{(\Composition{T}{S})}
    = \Composition{(\Composition{S_1}{T})}{S}
    = \Composition{(\Composition{S_2}{T})}{S}
    = \Composition{S_2}{(\Composition{T}{S})} = S_2$.

    \textbf{(12) $\Rightarrow$ (1)}:
    Suppose $T$ is epic. Pick a complement $W$ of $\Rng{T}$ in $V$, so
    $V = \Rng{T} \DirectSum W$. Let $\pi: V \to W$ be the projection
    ($\pi(r + w) = w$). For any $u \in U$: $\pi(T(u)) = 0$ since
    $T(u) \in \Rng{T}$. Thus $\Composition{\pi}{T} = \Composition{0}{T}$.
    By the epic property, $\pi = 0$, so $W = \{0\}$ and $\Rng{T} = V$.
\end{proof}

\begin{remark}[Split epimorphism and epimorphism, diagrammatically]
    Dually to the injective case, condition~(11), that $T$ is a \textbf{split epimorphism},
    asserts a right inverse (section) $S$ making the triangle commute:
    \begin{center}
    \begin{tikzcd}
        V \arrow[r, "S"] \arrow[dr, "\Identity{V}"'] & U \arrow[d, "T"] \\
        & V
    \end{tikzcd}
    \end{center}
    Condition~(12), that $T$ is an \textbf{epimorphism}, is right-cancellability of $T$ against
    any parallel pair $S_1, S_2: V \to W$:
    \begin{center}
    \begin{tikzcd}
        U \arrow[r, "T"] & V \arrow[r, shift left, "S_1"] \arrow[r, shift right, "S_2"'] & W
    \end{tikzcd}
    \qquad
    $\Composition{S_1}{T} = \Composition{S_2}{T} \;\implies\; S_1 = S_2$.
    \end{center}
    A section always makes $T$ epic; over a field the converse holds as well, which is the
    substance of the equivalence (11)$\Leftrightarrow$(12). The arrows are exactly those of
    the split-monomorphism diagram with their direction reversed --- the formal expression of
    injectivity/surjectivity duality.
\end{remark}

\newpage
\begin{theorem}[Characterizations of Bijectivity]\label{thm:bijectivity}
    For a linear map $T: U \to V$, the following are equivalent. Apart from the defining
    condition~(2), each is the conjunction of a characterization of injectivity
    (Theorem~\ref{thm:injectivity}) with the corresponding characterization of surjectivity
    (Theorem~\ref{thm:surjectivity}):
    \begin{enumerate}
        \item $\Ker{T} = \{0\}$ \textbf{and} $\Rng{T} = V$
        \item $T$ is bijective
        \item $\InverseImage{T}{\DirectImage{T}{W}} = W$ for every $W \Subspace[K] U$
              \textbf{and} $\DirectImage{T}{\InverseImage{T}{W'}} = W'$ for every
              $W' \Subspace[K] V$
        \item $W \mapsto \DirectImage{T}{W}$ is injective on the subspace lattice of $U$
              \textbf{and} $W' \mapsto \InverseImage{T}{W'}$ is injective on the subspace
              lattice of $V$ (equivalently, the two are mutually inverse lattice isomorphisms)
        \item $\Dim[K]{\DirectImage{T}{W}} = \Dim[K]{W}$ for every $W \Subspace[K] U$
              \textbf{and} $\Codim[K]{\InverseImage{T}{W'}}{U} = \Codim[K]{W'}{V}$ for every
              $W' \Subspace[K] V$
        \item $\Dim[K]{\DirectImage{T}{W}} = \Dim[K]{W}$ for every \emph{finite-dimensional}
              $W \Subspace[K] U$ \textbf{and} $\Codim[K]{\InverseImage{T}{W'}}{U} =
              \Codim[K]{W'}{V}$ for every \emph{finite-codimensional} $W' \Subspace[K] V$
        \item $\Dim[K]{\DirectImage{T}{W}} = \Dim[K]{W}$ for every \emph{one-dimensional}
              $W \Subspace[K] U$ \textbf{and} $\Codim[K]{\InverseImage{T}{W'}}{U} =
              \Codim[K]{W'}{V}$ for every \emph{hyperplane} $W' \Subspace[K] V$
        \item $T$ sends every basis of $U$ to a basis of $V$
        \item $T$ sends some basis of $U$ to a basis of $V$
        \item $T$ is an isomorphism: there is a \emph{unique} linear map $T^{-1}: V \to U$
              with $\Composition{T^{-1}}{T} = \Identity{U}$ and
              $\Composition{T}{T^{-1}} = \Identity{V}$
        \item $T$ is monic \textbf{and} epic
    \end{enumerate}
\end{theorem}

\begin{proof}
    Let $\mathrm{I}_k$ and $\mathrm{S}_k$ denote characterization~(k) of
    Theorems~\ref{thm:injectivity} and~\ref{thm:surjectivity}; within each of those theorems
    all characterizations are equivalent. We prove the chain
    (1)$\Rightarrow$(2)$\Rightarrow$(3)$\Rightarrow$(4)$\Rightarrow$(5)$\Rightarrow$(6)$\Rightarrow$(7)$\Rightarrow$(8)$\Rightarrow$(9)$\Rightarrow$(10)$\Rightarrow$(11)$\Rightarrow$(1),
    each step pairing an injectivity implication with the parallel surjectivity one.

    \textbf{(1) $\Rightarrow$ (2)}:
    $\Ker{T} = \{0\}$ gives injectivity ($\mathrm{I}_1 \Rightarrow \mathrm{I}_2$) and
    $\Rng{T} = V$ gives surjectivity ($\mathrm{S}_1 \Rightarrow \mathrm{S}_2$), so $T$ is
    bijective.

    \textbf{(2) $\Rightarrow$ (3)}:
    Injectivity yields $\InverseImage{T}{\DirectImage{T}{W}} = W$ for all $W \Subspace[K] U$
    ($\mathrm{I}_2 \Rightarrow \mathrm{I}_3$); surjectivity yields
    $\DirectImage{T}{\InverseImage{T}{W'}} = W'$ for all $W' \Subspace[K] V$
    ($\mathrm{S}_2 \Rightarrow \mathrm{S}_3$).

    \textbf{(3) $\Rightarrow$ (4)}:
    These identities make $W \mapsto \DirectImage{T}{W}$ and $W' \mapsto \InverseImage{T}{W'}$
    injective on the respective subspace lattices
    ($\mathrm{I}_3 \Rightarrow \mathrm{I}_4$, $\mathrm{S}_3 \Rightarrow \mathrm{S}_4$); being
    mutually inverse by~(3) and inclusion-preserving, they are then mutually inverse lattice
    isomorphisms.

    \textbf{(4) $\Rightarrow$ (5)}:
    Injectivity of $W \mapsto \DirectImage{T}{W}$ gives
    $\Dim[K]{\DirectImage{T}{W}} = \Dim[K]{W}$ for all $W$
    ($\mathrm{I}_4 \Rightarrow \mathrm{I}_5$); injectivity of $W' \mapsto \InverseImage{T}{W'}$
    gives the corresponding codimension equality ($\mathrm{S}_4 \Rightarrow \mathrm{S}_5$).

    \textbf{(5) $\Rightarrow$ (6)}:
    Restrict the dimension equality to finite-dimensional $W$ and the codimension equality to
    finite-codimensional $W'$ ($\mathrm{I}_5 \Rightarrow \mathrm{I}_6$,
    $\mathrm{S}_5 \Rightarrow \mathrm{S}_6$).

    \textbf{(6) $\Rightarrow$ (7)}:
    Restrict further to one-dimensional $W$ and to hyperplanes $W'$
    ($\mathrm{I}_6 \Rightarrow \mathrm{I}_7$, $\mathrm{S}_6 \Rightarrow \mathrm{S}_7$).

    \textbf{(7) $\Rightarrow$ (8)}:
    In Theorem~\ref{thm:injectivity}, $\mathrm{I}_7 \Rightarrow \mathrm{I}_9$: $T$ sends every
    basis to a linearly independent set. In Theorem~\ref{thm:surjectivity},
    $\mathrm{S}_7 \Rightarrow \mathrm{S}_9$: $T$ sends every basis to a generating system.
    Since a basis is an independent generating set, $T$ sends every basis to a basis.

    \textbf{(8) $\Rightarrow$ (9)}:
    Every basis is some basis.

    \textbf{(9) $\Rightarrow$ (10)}:
    Let $B$ be a basis of $U$ with $\DirectImage{T}{B}$ a basis of $V$. Since
    $\DirectImage{T}{B}$ is linearly independent, $\Ker{T} = \{0\}$
    ($\mathrm{I}_{10} \Rightarrow \mathrm{I}_1$); since it generates $V$, $\Rng{T} = V$
    ($\mathrm{S}_{10} \Rightarrow \mathrm{S}_1$). Thus $T$ is bijective, and its set-theoretic
    inverse $T^{-1}: V \to U$ is linear: for $v_1, v_2 \in V$ and $a, b \in K$, putting
    $u_i = T^{-1}(v_i)$ gives $T(a u_1 + b u_2) = a v_1 + b v_2$, so
    $T^{-1}(a v_1 + b v_2) = a u_1 + b u_2 = a\,T^{-1}(v_1) + b\,T^{-1}(v_2)$. It is the unique
    two-sided inverse: if also $\Composition{R}{T} = \Identity{U}$ and
    $\Composition{T}{R} = \Identity{V}$, then
    $R = \Composition{R}{(\Composition{T}{T^{-1}})}
    = \Composition{(\Composition{R}{T})}{T^{-1}} = T^{-1}$.

    \textbf{(10) $\Rightarrow$ (11)}:
    $T^{-1}$ is a left inverse, so $T$ is monic ($\mathrm{I}_{11} \Rightarrow \mathrm{I}_{12}$),
    and a right inverse, so $T$ is epic ($\mathrm{S}_{11} \Rightarrow \mathrm{S}_{12}$).

    \textbf{(11) $\Rightarrow$ (1)}:
    Monic gives $\Ker{T} = \{0\}$ ($\mathrm{I}_{12} \Rightarrow \mathrm{I}_1$) and epic gives
    $\Rng{T} = V$ ($\mathrm{S}_{12} \Rightarrow \mathrm{S}_1$).
\end{proof}

\begin{remark}[Isomorphism, diagrammatically]
    Bijectivity (condition~(10)) fuses a retraction and a section into a single two-sided
    inverse: an \textbf{isomorphism} is a map $T$ together with $T^{-1}$ making both round
    trips the identity,
    \begin{center}
    \begin{tikzcd}[column sep=large]
        U \arrow[r, shift left, "T"] & V \arrow[l, shift left, "T^{-1}"]
    \end{tikzcd}
    \qquad
    $\Composition{T^{-1}}{T} = \Identity{U}, \qquad \Composition{T}{T^{-1}} = \Identity{V}$.
    \end{center}
    Equivalently, $T$ is at once a split monomorphism and a split epimorphism, and its left and
    right inverses then necessarily coincide (the computation in step~(9)$\Rightarrow$(10)
    above).
\end{remark}

\begin{definition}[Inverse Linear Map]\label{def:inverse-map}
    Let $T: U \to V$ be a bijective linear map. Its \textbf{inverse linear map}
    $T^{-1}: V \to U$ is the unique linear map such that
    $\Composition{T^{-1}}{T} = \Identity{U}$ and $\Composition{T}{T^{-1}} = \Identity{V}$.
\end{definition}

\begin{definition}[Isomorphism]\label{def:isomorphism}
    A linear map $T: U \to V$ is an \textbf{isomorphism} if it has a two sided linear inverse.
    We write $U \cong V$ if there exists an isomorphism between $U$ and $V$.
\end{definition}

\begin{corollary}[Dimension and Isomorphism]\label{cor:dim-iso}
    Two vector spaces are isomorphic if and only if they have the same dimension.
    In particular, every $n$-dimensional $K$-vector space is isomorphic to $K^n$.
\end{corollary}

\begin{proof}
    ($\Rightarrow$)
    If $T: U \to V$ is an isomorphism and $B$ is a basis for $U$, then $\DirectImage{T}{B}$ is a basis for $V$
    (Theorem~\ref{thm:bijectivity}). Since $T$ is bijective, $\Card{B} = \Card{\DirectImage{T}{B}}$.
    Thus $\Dim[K]{U} = \Dim[K]{V}$.

    ($\Leftarrow$)
    If $\Dim[K]{U} = \Dim[K]{V}$, let $B_U$ and $B_V$ be bases for $U$ and $V$ respectively.
    There exists a bijection $f: B_U \to B_V$. By the universal property (Theorem~\ref{thm:vector-space-universal}),
    $f$ extends to a linear map $T: U \to V$.
    Since $T$ sends the basis $B_U$ to the basis $B_V$, $T$ is an isomorphism.
\end{proof}

\begin{theorem}[Equal Dimension Equivalence]\label{thm:equal-dim-equiv}
    Let $U$ and $V$ be \textbf{finite-dimensional} $K$-vector spaces with $\Dim[K]{U} = \Dim[K]{V}$,
    and let $T: U \to V$ be a linear map. Then the following are equivalent:
    \begin{enumerate}
        \item $T$ is injective
        \item $T$ is surjective
        \item $T$ is bijective
    \end{enumerate}
\end{theorem}

\begin{proof}
    Let $n = \Dim[K]{U} = \Dim[K]{V}$, and let $B = \{b_0, \ldots, b_{n-1}\}$ be a basis for $U$.

    \textbf{(1) $\Rightarrow$ (3)}: Suppose $T$ is injective.
    By Theorem~\ref{thm:injectivity}(8), $\DirectImage{T}{B}$ is linearly independent in $V$.
    Since $|\DirectImage{T}{B}| = n = \Dim[K]{V}$, it is a maximal linearly independent set,
    hence a basis for $V$ (Proposition~\ref{prop:basis-maxmin}). Thus $T$ is surjective, hence bijective.

    \textbf{(2) $\Rightarrow$ (3)}: Suppose $T$ is surjective.
    By Theorem~\ref{thm:surjectivity}(9), $\DirectImage{T}{B}$ generates $V$.
    Since a generating set has at least $\Dim[K]{V} = n$ elements and $|\DirectImage{T}{B}| \leq n$,
    we have $|\DirectImage{T}{B}| = n$, so it is a minimal generating set, hence a basis.
    Thus $\DirectImage{T}{B}$ is linearly independent, so $T$ is injective, hence bijective.

    \textbf{(3) $\Rightarrow$ (1), (2)}: Trivial.
\end{proof}

\begin{corollary}[Endomorphism Equivalence]\label{thm:endo-equiv}
    Let $V$ be a \textbf{finite-dimensional} $K$-vector space and $T: V \to V$ a linear endomorphism.
    Then the following are equivalent:
    \begin{enumerate}
        \item $T$ is injective
        \item $T$ is surjective
        \item $T$ is bijective
    \end{enumerate}
\end{corollary}

\begin{proof}
    Apply Theorem~\ref{thm:equal-dim-equiv} with $U = V$.
\end{proof}

\begin{example}[Counterexamples]\label{ex:endo-counterexamples}
    The theorem fails without the hypotheses:
    \begin{enumerate}
        \item \textbf{Infinite-dimensional, injective but not surjective}:
              The right-shift operator $R: K^{(\bb{N})} \to K^{(\bb{N})}$ defined by
              $R(a_0, a_1, a_2, \ldots) = (0, a_0, a_1, a_2, \ldots)$ is injective but not surjective
              ($(1, 0, 0, \ldots) \notin \Rng{R}$).

        \item \textbf{Infinite-dimensional, surjective but not injective}:
              The left-shift operator $L: K^{(\bb{N})} \to K^{(\bb{N})}$ defined by
              $L(a_0, a_1, a_2, \ldots) = (a_1, a_2, a_3, \ldots)$ is surjective but not injective
              ($L(1, 0, 0, \ldots) = (0, 0, \ldots) = L(0, 0, \ldots)$).

        \item \textbf{Unequal dimensions ($\Dim[K]{U} \neq \Dim[K]{V}$)}:
              What matters is not that domain and codomain differ as sets but that their
              dimensions do --- here $\Dim[K]{K^n} = n \neq n+1 = \Dim[K]{K^{n+1}}$.
              The inclusion $\iota: K^n \hookrightarrow K^{n+1}$, $\iota(x_0, \ldots, x_{n-1}) = (x_0, \ldots, x_{n-1}, 0)$
              is injective but not surjective.
              The projection $\pi: K^{n+1} \to K^n$, $\pi(x_0, \ldots, x_n) = (x_0, \ldots, x_{n-1})$
              is surjective but not injective.
    \end{enumerate}
\end{example}

\begin{theorem}[Natural Isomorphism ${\Hom[K]{K}{V}} \cong V$]\label{thm:hom-k-v-iso}
    Let $V$ be a $K$-vector space. The evaluation map
    \[
        \varepsilon: \Hom[K]{K}{V} \to V, \quad T \mapsto T(1_K)
    \]
    is a $K$-linear isomorphism.
\end{theorem}

\begin{proof}
    \textbf{Linearity}: For $T, S \in \Hom[K]{K}{V}$ and $\alpha, \beta \in K$:
    \[
        \varepsilon(\alpha T + \beta S)
        = (\alpha T + \beta S)(1_K)
        = \alpha T(1_K) + \beta S(1_K)
        = \alpha \varepsilon(T) + \beta \varepsilon(S)
    \]

    \textbf{Injectivity}: Suppose $\varepsilon(T) = T(1_K) = 0$.
    For any $\lambda \in K$:
    \[
        T(\lambda) = T(\lambda \cdot 1_K) = \lambda T(1_K) = \lambda \cdot 0 = 0
    \]
    Thus $T = 0$, so $\Ker{\varepsilon} = \{0\}$.

    \textbf{Surjectivity}: For any $v \in V$, define $T_v: K \to V$ by
    $T_v(\lambda) = \lambda v$. Then $T_v$ is linear:
    \[
        T_v(\alpha \lambda + \beta \mu)
        = (\alpha \lambda + \beta \mu) v
        = \alpha (\lambda v) + \beta (\mu v)
        = \alpha T_v(\lambda) + \beta T_v(\mu)
    \]
    and $\varepsilon(T_v) = T_v(1_K) = 1_K \cdot v = v$.
\end{proof}

\begin{remark}
    This isomorphism is \textbf{natural}: it does not depend on any choice of basis.
    Every linear map $T: K \to V$ is completely determined by $T(1_K)$,
    since $T(\lambda) = \lambda T(1_K)$ by linearity. The inverse
    $\varepsilon^{-1}: V \to \Hom[K]{K}{V}$ sends $v \mapsto T_v$
    where $T_v(\lambda) = \lambda v$.
\end{remark}

\begin{remark}[Function application is a special case of composition]
\label{rem:application-is-composition}
    The isomorphism $\varepsilon$ does more than count dimensions: it lets us regard each vector
    $v \in V$ as the linear map $T_v = \varepsilon^{-1}(v): K \to V$, $\lambda \mapsto \lambda v$ --- the
    \emph{name} of $v$. Under this identification, \textbf{function application becomes composition}. For
    any linear map $T: V \to W$ and any $v \in V$, the composite $\Composition{T}{T_v}: K \to W$ sends
    \[
        1_K \longmapsto T\bigl(T_v(1_K)\bigr) = T(v),
    \]
    so $\Composition{T}{T_v} = T_{T(v)}$, i.e.\ $\varepsilon(\Composition{T}{T_v}) = T(v)$. Evaluating
    $T$ at $v$ is the same as post-composing the name $T_v$ with $T$:
    \[
        \underbrace{T(v)}_{\text{application}}
        \quad\longleftrightarrow\quad
        \underbrace{\Composition{T}{T_v}}_{\text{composition}} .
    \]
    Thus application is not a separate operation sitting beside composition; it is exactly the case of
    composition in which the source is the ground field $K$. This is the seed of a later fact: once linear
    maps are represented by matrices, the action of a matrix on a column vector $v \mapsto Tv$ is precisely
    this composition, and so ``matrix times vector'' is a special case of matrix multiplication rather
    than a new primitive.
\end{remark}

\begin{theorem}[Coordinate Isomorphism]\label{thm:coord-iso}
    Let $\mathcal{B} = (b_i)_{i \in I}$ be an ordered basis for $V$. The coordinate map
    \[
        \Chart{\mathcal{B}}: V \to K^{(I)}, \quad v \mapsto \CoordVec[\mathcal{B}]{v}
    \]
    is a $K$-linear isomorphism.
\end{theorem}

\begin{proof}
    \textbf{Linearity}: For $v = \sum_{i \in I} a_i b_i$ and $w = \sum_{i \in I} c_i b_i$ (with finite supports), and $\alpha, \beta \in K$:
    \[
        \alpha v + \beta w = \sum_{i \in I} (\alpha a_i + \beta c_i) b_i
    \]
    so $\Chart{\mathcal{B}}(\alpha v + \beta w) = (\alpha a_i + \beta c_i)_{i \in I}
    = \alpha (a_i)_{i \in I} + \beta (c_i)_{i \in I}
    = \alpha \Chart{\mathcal{B}}(v) + \beta \Chart{\mathcal{B}}(w)$.

    \textbf{Injectivity}: If $\Chart{\mathcal{B}}(v) = 0$, then $v = \sum_{i \in I} 0 \cdot b_i = 0$.

    \textbf{Surjectivity}: For any $(a_i)_{i \in I} \in K^{(I)}$ (finitely supported),
    $v = \sum_{i \in I} a_i b_i \in V$ satisfies $\Chart{\mathcal{B}}(v) = (a_i)_{i \in I}$.
\end{proof}

\begin{corollary}[Coordinate Bijection]\label{cor:coord-bijection}
    The coordinate map $\Chart{\mathcal{B}}$ is a bijection between $V$ and $K^{(I)}$.
    In particular, $V \cong K^{(I)}$ as $K$-vector spaces.
\end{corollary}

\begin{theorem}[Ordered Bases as Isomorphisms]\label{thm:bases-as-isos}
    Let $V$ be a $K$-vector space and $I$ an index set. The assignment sending an ordered basis
    $\mathcal{B} = (b_i)_{i \in I}$ of $V$ to the unique $K$-linear map
    \[
        \Frame{\mathcal{B}} : K^{(I)} \to V, \qquad \Frame{\mathcal{B}}(\CanonicalVector{i}) = b_i
        \quad (i \in I),
    \]
    is a bijection
    \[
        \Set{\text{ordered bases of } V \text{ indexed by } I}
        \ \xrightarrow{\ \sim\ }\ \operatorname{Iso}_K\!\left(K^{(I)},\, V\right)
    \]
    onto the set of $K$-linear isomorphisms $K^{(I)} \to V$, with inverse $\psi \mapsto
    (\psi(\CanonicalVector{i}))_{i \in I}$ (the image of the canonical basis). Taking
    $\Card{I} = \Dim[K]{V}$, the ordered bases of $V$ are thus in canonical bijection with the
    isomorphisms $K^{(\Dim[K]{V})} \to V$; moreover $\Frame{\mathcal{B}} = \bigl(\Chart{\mathcal{B}}\bigr)^{-1}$
    is the inverse of the coordinate isomorphism of Theorem~\ref{thm:coord-iso}.
\end{theorem}

\begin{proof}
    \textbf{Well defined.} Given an ordered basis $\mathcal{B} = (b_i)_{i \in I}$, the canonical basis
    $(\CanonicalVector{i})_{i \in I}$ of $K^{(I)}$ determines a \emph{unique} linear map
    $\Frame{\mathcal{B}}$ with $\Frame{\mathcal{B}}(\CanonicalVector{i}) = b_i$ (a linear map is freely
    determined by prescribing arbitrary values on a basis). It sends the basis
    $(\CanonicalVector{i})_{i \in I}$ to the basis $\mathcal{B}$, hence is an isomorphism by the basis
    characterization of bijectivity (Theorem~\ref{thm:bijectivity}); thus
    $\Frame{\mathcal{B}} \in \operatorname{Iso}_K(K^{(I)}, V)$.

    \textbf{Inverse.} Write $\Psi(\mathcal{B}) = \Frame{\mathcal{B}}$ and
    $\Theta(\psi) = (\psi(\CanonicalVector{i}))_{i \in I}$. For an isomorphism $\psi$, the family
    $(\psi(\CanonicalVector{i}))_{i \in I}$ is an ordered basis of $V$, since an isomorphism carries a
    basis to a basis (Theorem~\ref{thm:bijectivity}); so $\Theta$ is well defined. Now
    $\Theta(\Psi(\mathcal{B})) = (\Frame{\mathcal{B}}(\CanonicalVector{i}))_{i \in I} = (b_i)_{i \in I}
    = \mathcal{B}$, and for any isomorphism $\psi$ the maps $\Psi(\Theta(\psi))$ and $\psi$ agree on the
    basis $(\CanonicalVector{i})_{i \in I}$, hence are equal. Therefore $\Psi$ and $\Theta$ are mutually
    inverse, and $\Psi$ is a bijection.

    Finally, $\Frame{\mathcal{B}}(\CanonicalVector{i}) = b_i$ gives
    $\Chart{\mathcal{B}}(b_i) = \CanonicalVector{i} = \Frame{\mathcal{B}}^{-1}(b_i)$, so
    $\Chart{\mathcal{B}} = \Frame{\mathcal{B}}^{-1}$.
\end{proof}

\begin{remark}[Bases are a choice of coordinates]
    Theorem~\ref{thm:bases-as-isos} is the precise sense in which ``choosing an ordered basis of $V$''
    and ``choosing an isomorphism $K^{(I)} \to V$'' are the same act: a basis \emph{is} an
    identification of $V$ with a standard coordinate space. This identification is the foundation of the
    next chapter, where it turns linear maps into matrices.
\end{remark}
